\section{VAE architecture for M1: end-to-end vs.\ split backbone+head}
\label{sec:arch}

\Cref{sec:use} sketches the M1 prior and loss abstractly, on top of whatever encoder/decoder
produces $\hat Z$ from $X$. \Cref{thm:ng} and the propositions built on it in this note are stated for
exactly one generative object: $X=g(Z)$, with $X$ the actual observed image and $g$ the true, unknown
mixing map. Only one of the two architectures compared in this section actually instantiates that
object. \Cref{sec:archA} (end-to-end) does. \Cref{sec:archB} (split, frozen backbone) does not: it
reconstructs a backbone-derived quantity, not $X$, and nothing in \cite{ng2025} says anything about the
resulting model (\cref{sec:archB:gap}). \Cref{sec:archB} is kept in this note anyway, but as what it
is -- an unproven speed heuristic, not a second way of applying the theory -- and this section is
written to say so plainly rather than to imply otherwise. \Cref{app:code} is used throughout as the
concrete description of what the reference implementation of Ng et al. \cite{ngimpl} already provides
and what it is missing.

\subsection[Option A: end-to-end]{Option A: end-to-end -- the only option that implements \cref{def:m1} and \cref{thm:ng} as stated}
\label{sec:archA}
The backbone and the VAE head are one model, trained jointly on raw $224\times224$ images: a
ResNet-50 feeds \texttt{encoder\_mu}/\texttt{encoder\_logvar}, and every term of the M1 loss
(reconstruction, KL, sparsity; \cref{app:code} on \texttt{Trainer.step}) backpropagates through the
head \emph{and} through the backbone. This is the only one of the two options in which the object
being reconstructed -- and hence the object the identifiability argument is actually about -- is the
true image $x$: the composite (backbone $+$ head) plays the role of the encoder, its decoder plays the
role of $g$, and the loss compares $\hat x$ against the actual $x$, exactly $X=g(Z)$ of \cref{def:m1},
with no intermediate object substituted in. \Cref{fig:optA} shows the flow.

\begin{figure}[H]
\centering
\begin{tikzpicture}[
  >=Stealth,
  font=\small,
  procnode/.style={draw,rounded corners,align=center,inner sep=1.8mm,minimum height=8mm},
  smallnode/.style={draw,rounded corners,align=center,inner sep=1.3mm,font=\footnotesize,minimum height=7mm},
  obsnode/.style={circle,draw,fill=gray!25,minimum size=6.5mm,inner sep=1pt}
]

\node[procnode] (img) at (0.8,3) {224$\times$224\\image $x$};
\node[procnode] (bb) at (3.5,0 |- img.east) {%
  \begin{tabular}{@{}c@{}} ResNet-50 \\ backbone \\ (trainable) \end{tabular}};
\node[smallnode] (enc) at (6.6,0 |- bb.east) {%
  \begin{tabular}{@{}c@{}} encoder\_mu \\ encoder\_logvar \end{tabular}};
\node[procnode] (z) at (9.3,0 |- enc.east) {sample $z$};
\node[procnode] (dec) at (0.3,4.4 -| z.north) {decoder};
%\node[procnode] (xhat) at (12.1,0 |- dec.east) {$\hat x$};

\node[smallnode] (meannet) at (5.6,0.9) {%
  \begin{tabular}{@{}c@{}} mean\_net \\ (adj-masked $z$) \end{tabular}};
\node[obsnode] (u) at (8.6,0 |- meannet.east) {$u$};
\node[smallnode] (logvarnet) at (10.6,0 |- u.east) {%
  \begin{tabular}{@{}c@{}} log\_var\_net \\ (domain $u$) \end{tabular}};

\node[procnode] (loss) at (9.0,-1.3) {%
  \begin{tabular}{@{}c@{}} loss \\ recon$(x,\hat x)$ + KL + sparsity \end{tabular}};

\draw[->] (img) -- node[above] {$x$} (bb);
\draw[->] (bb) -- node[above] {$\tilde x$} (enc);
\draw[->] (enc) -- (z);
\draw[->] (z) -- node[right] {$\hat z$} (dec);
%\draw[->] (dec) -- (xhat);
\draw[->] (z) -- (meannet);
\draw[->] (u) -- (logvarnet);
\draw[->] (meannet) -- (loss);
\draw[->] (logvarnet) -- (loss);

\coordinate (dropX1) at (12.1,0 |- dec.east);
\coordinate (dropX2) at (dropX1 |- loss);

\draw[->]
  (dec) -- (dropX1)
  -- node[left,pos=.4] {$\hat x$}
  (dropX2)
  -- (loss);

\coordinate (dropA1) at (loss.south |- 0,-2.5);
\coordinate (dropA2) at (dropA1.west -| 3.5,0);
\draw[->,dashed,red] (loss.south) -- (dropA1) -- (dropA2) -- (bb.south);
\node[font=\scriptsize,align=center] at (6.3,-2.8)
  {gradients: all parameters, including the backbone};

\coordinate (dropI1) at (img.south |- loss.west);
\draw[->]
  (img.south)
    -- node[right,midway] {$x$}
  (dropI1)
    -- (loss.west);

\end{tikzpicture}
\caption{\textbf{\nameref{sec:archA}} Every training step runs the full backbone forward, and the
M1 loss (dashed red) backpropagates through the VAE head and through the backbone. The reconstruction
term compares $\hat x$ against the true image $x$ -- the object \cref{def:m1,thm:ng} are stated about.}
\label{fig:optA}
\end{figure}
\newpage
\subsection[Option B: split]{Option B: split (frozen backbone, cached features) -- a heuristic, not a
second implementation of the theory}
\label{sec:archB}
Three separate stages: one to adapt the backbone, one (also run only once) to extract and cache
features, and one -- the actual M1 estimation -- that is trained and re-trained many times as the
prior/graph design is iterated. Write $x$ for the true image and $\tilde x:=\mathrm{BB}(x)$ for the
backbone's output, cached at stage 2 (the $\sim\!3000$-dim embedding described below). This notation
matters: as \cref{sec:archB:gap} makes precise, stage 3 never sees $x$ at all, only $\tilde x$, and
everything below should be read with that substitution in view.
\begin{enumerate}[label=(\arabic*),leftmargin=2em]
\item \textbf{Backbone adaptation (once).} Start from an ImageNet-pretrained ResNet-50. Adjust only its
      batch-norm (BN) layers to TerraIncognita (TI); every other weight stays frozen. 
      Concretely: run the network in
      \texttt{train()} mode (so BN layers update from batch statistics) but with gradients disabled,
      over batches pooled across the three training locations L38/L43/L46 (never L100), with BN 
      momentum set to \texttt{None} so each layer's running mean/variance is a cumulative 
      (equally-weighted) average of the per-batch statistics over the whole pass, rather than an 
      exponential moving average dominated by recent batches.
\item \textbf{Feature extraction (once).} Run the BN-adjusted, still-frozen backbone $\mathrm{BB}$ in
      \texttt{eval()} mode once over every image $x$ -- train and test, L100 included -- and cache
      $\tilde x=\mathrm{BB}(x)$ (per the settled pipeline: layer3 pooled concatenated with the final
      pooled layer, plus a brightness scalar) together with species/location/time-of-day/image-id in
      one table. This pass never touches a gradient and never needs to be repeated once the backbone
      from step (1) is fixed.
      \begin{itemize}
      \item \texttt{layer3} refers to ResNet-50's third residual stage 
      For a 224$\times$224 input:
      \begin{itemize} 
      \item \texttt{layer3} output: 1024 channels, 14$\times$14 spatial
	      $\rightarrow$ global-average-pooled to a 1024-dim vector
      \item \texttt{layer4} output (the ``final" stage before the classifier): 2048 channels, 7$\times$7 
	      spatial $\rightarrow$ pooled to 2048-dim -- this is the standard ``ResNet-50 feature vector"
	  \end{itemize}      
      Concatenating the two gives $\sim$3072 dims total, i.e.\ $\tilde x\in\R^{3072}$. The reason to
      include \texttt{layer3} and not just the usual final-layer 2048-dim vector: \texttt{layer4} 
      sits right before the (originally) 1000-way ImageNet classifier, so it's the most abstracted, 
      most class- \linebreak
      discriminative, and the stage most likely to have had low-level nuisance 
      information (illumination, color balance, camera signature) trained \textit{out} of it. 
      \texttt{layer3} 
      is one stage earlier and less abstracted, so it retains more texture/color/illumination 
      detail -- the kind of signal the M1 style block needs. Concatenating both gives one embedding 
      that's more likely to carry both content-relevant and style-relevant information, rather than 
      betting everything on whichever the last layer kept.
      \item \textit{Brightness} scalar: mean luminance of the image, computed as a plain scalar, not 
      from the network at all. Concretely:
      \begin{itemize}
      \item convert the image to luminance via a standard weighted sum, 
       		e.g. $Y = 0.299R + 0.587G + 0.114B$ (or BT.709 weights $0.2126/0.7152/0.0722$ -- either is 
       		fine, just stay consistent) 
      \item average $Y$ over all pixels $\rightarrow$ one number per image
      \end{itemize}
      Important detail: compute this on the raw resized image (pixel values in their normal range, e.g. 
      $[0,255]$ or $[0,1]$), before ImageNet mean/std normalization -- normalizing per-channel with 
      ImageNet's statistics would distort the brightness value in a way tied to ImageNet's own color 
      statistics, not the image's actual brightness. Same measure must be used at train and test time 
      (including on L100), since the note later proposes brightness as a TOD proxy -- consistency there 
      matters more than the exact formula chosen.
      \end{itemize}
\item \textbf{M1 estimation (iterated many times).} The M1 VAE (the small MLPs of \cref{app:code},
      modified per \cref{sec:use}) is trained on the cached tensor dataset $\tilde x$ from step (2);
      only those MLPs ever receive a gradient, and none of it touches the backbone. Its encoder maps
      $\tilde x\mapsto\hat Z$, and its decoder maps $\hat Z\mapsto\hat{\tilde x}$ -- a reconstruction of
      the cached embedding, not of the image. The reconstruction term of the loss is therefore
      $\mathrm{recon}(\tilde x,\hat{\tilde x})$, not $\mathrm{recon}(x,\hat x)$; \cref{fig:optB} draws
      this explicitly. This is the step that is actually redesigned and re-run repeatedly --
      content/style $u$-dependence, ANM vs.\ HNM, the sink-freezing rule -- while steps (1) and (2)
      stay fixed.
\end{enumerate}
\Cref{fig:optB} shows the flow.

\begin{figure}[H]
\centering
\begin{tikzpicture}[
  >=Stealth,
  font=\small,
  procnode/.style={draw,rounded corners,align=center,inner sep=1.8mm,minimum height=8mm},
  smallnode/.style={draw,rounded corners,align=center,inner sep=1.5mm,font=\footnotesize,minimum height=7mm},
  stagebox/.style={draw,dashed,rounded corners,inner sep=4mm}
]

% ---- stage 1: backbone adaptation ----
\node[procnode] (imnet) at (0,3.2) {%
  \begin{tabular}{@{}c@{}} ImageNet-\\pretrained \\ ResNet-50 \end{tabular}};
\node[procnode] (bnrecal) at (3.6,0 |- imnet.east) {%
  \begin{tabular}{@{}c@{}} BN-only recalibration \\ on L38/L43/L46 \\ (no gradient) \end{tabular}};
\node[procnode] (frozen) at (8.0,0 |- bnrecal.east) {%
  \begin{tabular}{@{}c@{}} frozen backbone $\mathrm{BB}$ \\ (all weights fixed) \end{tabular}};

\draw[->] (imnet) -- node[above] {$x$} (bnrecal);
\draw[->] (bnrecal) -- (frozen);

\begin{scope}[on background layer]
  \node[stagebox,fit=(imnet)(bnrecal)(frozen),inner sep=2mm] (stage1) {};
\end{scope}
\node[font=\small,above=1mm of stage1.north west,anchor=south west] {stage 1: backbone adaptation (once)};

% ---- stage 2: feature extraction (once) ----
\node[procnode] (allimg) at (0,0) {%
  \begin{tabular}{@{}c@{}} every image $x$ \\ (train \emph{and} L100) \end{tabular}};
\node[procnode] (cache) at (4.4,0 |- allimg.east) {%
  \begin{tabular}{@{}c@{}} cached table: $\tilde x=\mathrm{BB}(x)$ \\ + species/TOD/\\location/id \end{tabular}};

\draw[->] (frozen) -- ++(0,-1.6) -| (allimg.north);
\draw[->] (allimg) -- (cache);

\begin{scope}[on background layer]
  \node[stagebox,fit=(allimg)(cache),inner sep=2mm] (stage2) {};
\end{scope}
\node[font=\small,below=1mm of stage2.south west,anchor=north west] {stage 2: feature extraction (once)};

% ---- stage 3: M1 estimation (iterated many times) ----
\node[smallnode,minimum width=3.2cm] (m1enc) at (8.6,-2.9) {%
  \begin{tabular}{@{}c@{}} M1 VAE encoder: \\ encoder\_mu/logvar \\ (input: $\tilde x$) \end{tabular}};
\node[procnode] (zhat) at (12.0,0 |- m1enc.east) {sample $\hat z$};
\node[smallnode,minimum width=3.2cm] (m1dec) at (m1enc.south |- 0,-4.7) {%
  \begin{tabular}{@{}c@{}} M1 VAE decoder \\ + mean\_net/log\_var\_net \end{tabular}};
%\node[procnode] (xtilhat) at (zhat.south |- m1dec.east) {$\hat{\tilde x}$};
\node[procnode] (loss2) at (zhat.south |- 0,-6.4) {%
  \begin{tabular}{@{}c@{}} loss \\ recon$(\tilde x,\hat{\tilde x})$ + KL + sparsity \end{tabular}};

\coordinate (dropXt1) at (cache.east -| m1enc.north);
\draw[->]
  (cache.east) -- (dropXt1)
  -- node[right,pos=.5] {$\tilde x$}
  (m1enc.north);
\draw[->] (cache.east) -| (m1enc.north);
\draw[->] (m1enc) -- (zhat);
\draw[->] (zhat) -- node[right=4pt,pos=.5] {$\hat z$} (m1dec);
%\draw[->] (m1dec) -- (xtilhat);
%\draw[->] (xtilhat) -- (loss2);

\coordinate (dropXht1) at (zhat.south |- m1dec.east);
\draw[->] (m1dec.east) -- (dropXht1) -- node[right,pos=.5] {$\hat{\tilde{x}}$} (loss2);

\coordinate (dropT1) at (cache.south |- loss2.west);
\draw[->] (cache.south) -- (dropT1) -- (loss2.west);
\node[font=\scriptsize,align=center] at (2.6,-4.4)
  {recon target is $\tilde x$, \textbf{never} $x$};

\coordinate (dropB1) at (loss2.south |- 0,-7.6);
\coordinate (dropB2) at (m1dec.south |- dropB1.west);
\coordinate (dropB3) at (6.2,0 |- dropB1.west);
\draw[->,dashed,red] (loss2.south) -- (dropB1) -- (dropB2) -- (m1dec.south);
\draw[->,dashed,red] (dropB2) -- (dropB3) |- (m1enc.west);
\node[font=\scriptsize,align=center] (fbBlabel) at (10.3,-7.9)
  {gradients: MLPs only, never reach $\mathrm{BB}$};

\begin{scope}[on background layer]
  \node[
    stagebox,
    fit={
      (m1enc)(m1dec)(zhat)(loss2)(fbBlabel)
      ([xshift=-5mm]m1enc.west)
      ([xshift=-5mm]m1dec.west)
    },
    inner sep=2mm
  ] (stage3) {};
\end{scope}
\node[font=\small,below=1mm of stage3.south west,anchor=north west] {stage 3: M1 estimation (iterated many times)};

\end{tikzpicture}
\caption{\textbf{\nameref{sec:archB}} Stage 1 happens once and touches only BN statistics. Stage 2's
extraction pass also happens once and produces the cached table of $\tilde x=\mathrm{BB}(x)$. Stage 3
never sees $x$: its encoder takes $\tilde x$, its decoder reconstructs $\hat{\tilde x}$, and the
reconstruction term of the loss compares $\tilde x$ against $\hat{\tilde x}$, not $x$ against
$\hat x$ -- unlike \cref{fig:optA}. Gradients (dashed red) reach only the stage-3 MLPs and never the
backbone.}
\label{fig:optB}
\end{figure}

\paragraph{Why BN-only, not full fine-tuning.}
BN statistics (per-channel mean and variance) encode low-level image statistics -- illumination,
color balance, sensor response -- that differ between ImageNet and camera-trap imagery.
Recalibrating them closes much of that gap cheaply, without touching the learned convolutional
filters or anything built on top of them, needs no labels, and needs only one forward pass. It also
cannot overfit to three locations the way updating the convolutional weights (or the whole network)
could. This last point matters specifically for M1: a supervised fine-tuning objective (species
cross-entropy) has no reason to preserve location-relevant information and is, if anything, pushed to
discard it -- which would remove exactly the signal the style block $Z_s$ needs. BN-only adaptation
changes global per-channel statistics, not which features are computed, so it is far less likely to
erase that signal.

\paragraph{What reconstructing $\tilde x$ instead of $x$ actually means.}
\label{sec:archB:gap}
\Cref{def:m1,thm:ng} are stated for $X=g(Z)$ with $X$ the true observation and $g$ a $\mathrm{C}^2$
diffeomorphism onto its image (\cref{sec:ng}). Stage 3, as just described, does not fit that model. It
fits a model of $\tilde X=\mathrm{BB}(g^(Z))$ -- the composite of the true mixing $g$ with the 
backbone $\mathrm{BB}$ -- and \cite{ng2025} says nothing about this composite: $\mathrm{BB}$ is not part 
of the generative process \cref{def:m1} describes, and nothing constrains $\mathrm{BB}\circ g$ 
to be injective, differentiable in the required sense, or otherwise well-behaved, even when $g$ itself 
is.
Quite the opposite is likely: $\mathrm{BB}$ is trained (\cref{sec:pretrain}) on an objective with no
term rewarding invertibility, and both pretraining families discussed there are pushed, to different
degrees, toward \emph{discarding} exactly the nuisance/style variation that would need to survive for
$\mathrm{BB}\circ g$ to remain injective in the directions $Z_s$ occupies. \Cref{app:probe} tests
whether species and location remain \emph{linearly readable} in $\tilde x$ -- a much weaker property
than the injectivity the theorem needs, and one that can hold even when $\mathrm{BB}\circ g$ 
collapses many distinct $(Z_c,Z_s)$ onto embeddings a linear classifier can still separate but a decoder 
cannot invert.

One more point worth being explicit about, since it is easy to read the reference code the wrong way:
\cite{ng2025}'s own paper states the reconstruction loss results from the assumed 
decoder likelihood $p^{(u)}(X \mid \hat Z)$ (Gaussian, pre-specified variance), which 
doesn't hurt the identifiability argument but merely limits the scope. \nameref{sec:archB} breaks this: 
it takes that somewhat limited scope and turns it into an non-principled substitution ($\tilde x$ for 
$x$). 
\nameref{sec:archB} should therefore be read as exactly what it is: a speed-motivated heuristic that may 
or may not preserve enough of $Z$'s structure in $\tilde x$ for stage 3 to recover something useful, not 
a second, equally valid way of instantiating \cref{def:m1,thm:ng}. Whether it works is an empirical 
question this note does not resolve; \nameref{sec:archA} remains the only construction its own 
theoretical machinery actually describes.

\subsection{Pros and cons}

\paragraph{\nameref{sec:archA}}
\emph{Pros:} this is the only construction that actually instantiates $X=g(Z)$ as \cref{def:m1,thm:ng}
state it; the backbone can adapt jointly with the head, in principle giving the encoder more freedom to
place exactly the information the content/style split wants at the bottleneck; one unified model, no
separate caching step, no intermediate object standing in for $x$.
\emph{Cons:} every iteration runs a full ResNet-50 forward and backward pass on raw images -- hours,
not minutes, in your own experience -- while the M1 estimator itself is still being iterated (the
five open items closing \cref{app:code}: $u$-dependent offset/mean, per-coordinate mechanisms, the
sink-freezing test, the optimizer overlap, the dead code), so every one of those iterations would
carry the full backbone cost; the loss (\mbox{reconstruction + KL + sparsity}) has no term that 
specifically rewards keeping location information, so end-to-end gradient pressure can push the backbone 
to discard exactly the style-relevant cues $Z_s$ is supposed to capture; two learning problems with
different dynamics -- representation learning and causal/structural discovery (Adam on most
parameters, a separate SGD gate optimizer on the adjacency matrix, \cref{app:code}) -- are coupled
through shared backbone gradients, harder to diagnose than training the small MLPs alone; only three
training locations, so fine-tuning the whole backbone risks overfitting its features to those three
sites rather than to the location-invariant/location-varying split M1 assumes; and even here, nothing
guarantees the backbone itself is injective (any convolutional network with pooling generically is
not) -- end-to-end training only supplies gradient \emph{pressure} toward retaining what reconstruction
needs, not a proof that it does.

\paragraph{\nameref{sec:archB}}
\emph{Pros:} iteration is minutes rather than hours, which matters given how much of the M1 estimator
is still design rather than implementation. Extraction is a single pass per image, cached once;
everything downstream is plain backprop on $\sim\!3000$-dim vectors, cheap to re-run many times over.
\Cref{app:probe}'s linear-probe diagnostics give a cheap, checkable precondition test -- whether
species/location information survives at all in $\tilde x$ -- before committing to this heuristic for a
given backbone; passing that check narrows candidates but, as \cref{sec:archB:gap} explains, does not
establish that stage 3's fitted model bears any proven relationship to the true $(Z_c,Z_s,\Gc)$.
\emph{Cons:} \cref{sec:archB:gap} is the central one and subsumes the more specific concerns below --
this option does not implement \cref{def:m1,thm:ng} at all, and no identifiability result in this note
applies to it. More specifically: a hard ceiling set by the backbone -- whatever style/location
information is not present in $\tilde x$, even after BN-only recalibration, cannot be recovered by any
downstream causal method; $\tilde x$ is not shaped by the M1 reconstruction objective itself, so subtle
style cues a jointly-trained encoder might have preserved for reconstruction's sake could be missing;
changing the layer(s), pooling, or backbone means re-extracting and re-caching the whole table (cheap,
but a step that end-to-end training does not need).

\subsection{Recommendation}
\label{sec:rec}
Use \nameref{sec:archB} first, but as a heuristic, not as a theory-matching choice:
\cref{sec:archB:gap} means nothing in this note guarantees it recovers anything. The reasons to start
there are practical only. First, iteration speed: \cref{app:code} closes with five concrete open
changes to the reference implementation, each of which still needs to be tried and likely revised,
which is only practical at minutes per run rather than hours. Second, it decouples backbone selection
(checkable once, cheaply, via \cref{app:probe}) from estimator design (iterated many times), so a bad
backbone choice does not have to be discovered by watching the M1 estimator itself fail to converge.
Third, \nameref{sec:archA} does not automatically avoid the injectivity problem either -- any
convolutional backbone with pooling is generically non-injective regardless of how it is trained -- so
paying its much higher cost does not, by itself, buy a proof that the resulting model matches
\cref{def:m1}; it only removes the specific, additional substitution of $\tilde x$ for $x$ that
\cref{sec:archB:gap} describes, restoring $x$ as the reconstruction target the theorem is actually
about.

Given this, treat \nameref{sec:archB} as an exploratory, cheap-to-iterate first pass -- useful for
developing and debugging the prior/graph design of \cref{sec:use} -- and \nameref{sec:archA} as the one
construction whose relationship to \cref{def:m1,thm:ng} this note can actually state precisely: to be
used once the design stabilizes enough to justify its cost, and to check whether conclusions reached
under \nameref{sec:archB} survive when $x$, rather than $\tilde x$, is reconstructed.

\subsection{Pretraining objective for the backbone: supervised vs.\ self-supervised}
\label{sec:pretrain}
What matters for M1 is not raw discriminative accuracy on ImageNet classes but whether the embedding
still contains \emph{both} kinds of signal the split needs to separate afterward: content (species,
time of day) and style (location/camera cues). This choice matters most for \nameref{sec:archB}, where
whatever the backbone discards is discarded permanently -- stage 3 never touches $\mathrm{BB}$ again.
Under \nameref{sec:archA} the same choice still matters as the initialization the end-to-end
fine-tuning starts from, though there gradient pressure from the M1 loss can, in principle but without
guarantee, partially recover discarded signal. If either kind of signal is missing after pretraining
(and, for \nameref{sec:archB}, after BN-only recalibration), no downstream causal machinery in
\cref{app:code} can recover it -- this is exactly the feature-extraction ceiling noted as a con of
\nameref{sec:archB} above (\cref{sec:archB:gap}).

\paragraph{Supervised ImageNet pretraining.}
The objective is cross-entropy on object categories that do not correlate with camera exposure,
background statistics, or viewpoint, so nothing in the loss rewards keeping that information;
supervised features are, generically, pushed toward invariance to exactly this kind of within-class
nuisance variation. For M1 this is close to the worst case among the candidates: the pretraining
objective actively works against keeping the style signal the style block needs.

\paragraph{Self-supervised: SimCLR, MoCo v2, MoCo v3, SimSiam.}
All four train invariance to a shared, standard augmentation recipe -- random-resized crop, color
jitter, occasional grayscale, Gaussian blur, horizontal flip -- rather than invariance to a fixed
label set. What differs between them is the loss and the negative-pair mechanics (InfoNCE with a
large batch for SimCLR; a momentum encoder plus a memory queue for MoCo v2; a momentum encoder plus a
predictor without a memory queue for MoCo v3; a predictor with stop-gradient and no negatives at all
for SimSiam), which mainly affects training stability and representation quality for the downstream
discriminative task, not \emph{which} low-level factors get suppressed. So there is no a priori
reason to expect any one of the four to retain systematically more or less location/style information
than the others. Two caveats, both empirical rather than settled by the choice of method: (a) color
jitter and blur do push toward suppressing some of the brightness/exposure cues that matter for
time-of-day/style in TI, less aggressively than a labeled cross-entropy objective but not zero; (b)
none of the four ever see location or time-of-day labels during pretraining, so whatever style signal
survives is incidental to the augmentation recipe, not something the objective was built to preserve.

\paragraph{Recommendation.}
Prefer self-supervised pretraining over supervised, generically, since supervised is the one option
whose objective actively works against what M1 needs. Among SimCLR, MoCo v2, MoCo v3 and SimSiam
there is no first-principles reason to prefer one for style preservation; choose empirically, using
the linear-probe step of \cref{app:probe} (species accuracy plus within-species location accuracy on
cached features) over whichever of the four you already have trained, and keep the supervised
ResNet-50 only as a baseline/comparison point rather than a primary candidate.

\subsection{Verification of Recommendations}
\label{sec:verify}
Before investing in \nameref{sec:archB} for a given backbone, run the test outlined in \cref{app:probe}
to compare candidates. Per \cref{sec:archB:gap}, a pass there rules out only the worst case -- a
backbone that has discarded species or location information outright -- it does not establish that the
backbone, or its composition $\mathrm{BB}\circ g$ with the true mixing, is close to the injective map \cref{thm:ng} requires, and no diagnostic in this note checks that directly. \Cref{sec:archA} needs no equivalent gate, since it does not introduce $\tilde x$ or $\mathrm{BB}$ as a separate, unexamined step in the first place.
